\subsection*{Source position 0}
Finding out actual items bought via credit card issuer and not the store receipt?

\paragraph{Original; $p=3.02e-14$.} Credit card issuers do not automatically provide a direct list of every specific item purchased by the consumer. This information is strictly confidential because the actual purchase data resides with the merchant and the payment network, not the bank. Consequently, the card issuer can only tell you the total transaction amount, the date, and often the merchant's general category, but they cannot disclose the individual goods or services involved in the transaction. To find out what was actually bought, you must review your own digital or physical receipts or statements provided by the retail stores themselves, as the financial institution lacks visibility into the internal inventory or specific line-item details of the sale.


\paragraph{Pre / smoothing; quality passes; $p=5.42e-01$.} Credit card issuers cannot provide a comprehensive, itemized list of every purchase made by a consumer because such data is confidential and remains the property of the merchant and payment network rather than the bank. This limitation means that the issuer can only report the overall transaction total, the date of the transaction, and often a broad category, but is unable to reveal the specific goods or services purchased. To ascertain the exact nature of the items acquired, customers must consult their own retail receipts or statements, which serve as the only source available since the financial institution does not have access to the store's internal inventory or line-item breakdowns of the sale.

\paragraph{Conditional / smoothing; quality passes; $p=3.93e-03$.} Because credit card issuers are bound by strict confidentiality, they cannot furnish a comprehensive list of every transactional detail, such as the specific products acquired, since the primary seller and the payment processor hold the detailed purchase records that the bank does not. As a result, the financial institution is permitted to share only the aggregate sum of the charge, the transaction date, and the typical merchant classification, yet it remains unable to reveal the particular goods or services actually bought. To ascertain the exact contents of a purchase, an individual must rely on the receipts or statements issued directly by the retail establishments, as the credit card agency lacks access to the internal inventory of the merchants or the granular line-item specifics of the sale.

\subsection*{Source position 8}
Why do we spell the country as ' Philippines ' but the language as ' Filipino ' ? So why are they spelt differently ? ' Filipino ' also describes the people , culture , cuisine etc . Please explain like I'm five.

\paragraph{Original; $p=2.22e-11$.} We spell the country "Philippines" because that specific long word represents the entire nation, much like how the United States is a big group of many states. We spell the language "Filipino" simply because the shorter word "Phil" comes first; our brains naturally add the "i" to make it sound right and easier to say. Since the language, people, and food all speak this shorter "Phil" sound, we use "Filipino" for them. The long version stays for the land itself, while the shorter version grows into the culture and the spoken words we love.


\paragraph{Pre / smoothing; quality passes; $p=1.31e-01$.} We pronounce the nation "Philippines" as a single continuous name for the whole country, treating it like the United States, which stands for a collection of many states; in contrast, we call the language "Filipino" because the shorter name "Phil" begins first, and our innate tendency to add an "i" makes it a naturally pronounceable word. Since these phonetics are shared across the people, food, and the land, we designate the shortened version for the commoners, acknowledging that the original, lengthy spelling remains for the territory itself while the shorter form has evolved into the culture and everyday speech.

\paragraph{Conditional / smoothing; quality passes; $p=2.32e-06$.} Philippines is the full name of the country because the lengthy spelling captures the entire nation's scope, just as the United States aggregates many states. In contrast, the word Filipino itself is short, and because "Phil" already begins the name, our natural speech tendency inserts an "i" afterward to create a sound that is both easier to say and feels more natural. This shorter form, "Phil," becomes the foundation for the language, the people, and the cuisine, which all rely on that specific pronunciation. Consequently, the extensive original name remains reserved for the territory, while the reduced name expands to encompass the entire culture and the language we cherish.

\subsection*{Source position 16}
Please explain what is "Eric Schmidt"

\paragraph{Original; $p=8.09e-18$.} Eric Schmidt was a prominent figure in the technology industry, best known for serving as the founder and chief executive officer of Google from 2001 to 2011. His leadership strategy emphasized consumer innovation and expanding Google's global market presence while keeping the company focused on its core search business. He later co-founded Threadless, a groundbreaking graphic T-shirt company, before returning to serve as the chief executive officer of Altaba, a special-purpose acquisition company created to buy Intel's enterprise business. Ultimately, his career defined modern corporate strategy in the tech sector.


\paragraph{Pre / smoothing; quality passes; $p=3.41e-01$.} Eric Schmidt established himself as a key leader in the digital technology sector, notably serving as the co-founder and chief executive of Google between 2001 and 2011, where his approach prioritized consumer-driven innovation, expanded the firm's worldwide market reach while maintaining focus on search as its central offering. Subsequently, he co-established Threadless, an innovative graphic T-shirt brand, and then returned to lead Altaba as its chief executive officer; Altaba was a special-purpose acquisition vehicle designed to purchase Intel's enterprise division. In summary, Schmidt's professional trajectory fundamentally shaped the corporate strategies within the modern technology industry.

\paragraph{Conditional / smoothing; quality passes; $p=2.16e-02$.} Eric Schmidt was widely recognized as a leading professional in the technology sector, primarily for establishing Google and holding the role of chief executive officer between 2001 and 2011. During his tenure, he championed a management approach that prioritized consumer-driven innovation and the expansion of Google's worldwide market share, all while maintaining the organization's concentration on its core search operations. Following that period, he co-founded Threadless, an innovative graphic T-shirt manufacturer, and subsequently returned to lead Altaba as the chief executive officer, a special-purpose acquisition firm designed to acquire Intel's enterprise division. His professional path ultimately established him as a defining force in modern corporate strategy within the technology industry.

\subsection*{Source position 24}
Lend money at a rate linked to the prime rate

\paragraph{Original; $p=6.20e-08$.} Offering loans at a rate tied to the prime rate provides borrowers with a competitive advantage while maintaining institutional flexibility for lenders. As the prime rate fluctuates based on the federal funds rate, borrowers benefit from immediate, transparent pricing that reflects current market conditions without waiting for contractual resets. This approach simplifies communication and reduces administrative overhead compared to fixed rates, allowing financial institutions to attract price-sensitive customers effectively. Ultimately, this structure aligns borrowing costs with monetary policy shifts, ensuring that loan terms remain dynamically responsive to broader economic trends.


\paragraph{Pre / smoothing; quality passes; $p=8.44e-01$.} By linking loan interest to the prime rate, this mechanism offers lenders a significant degree of flexibility while delivering a competitive edge to borrowers who secure lower, more predictable rates. Since the prime rate is determined directly by the federal funds rate, it ensures that borrowers receive accurate and visible pricing that mirrors today's market reality in real time, eliminating the need to wait for contractual adjustments to apply. This dynamic approach streamlines customer outreach and lowers operational costs compared to fixed-rate options, enabling banks to successfully pull in clients who prioritize affordability. Consequently, the loan terms remain in constant harmony with changes in monetary policy, guaranteeing that the cost of borrowing adjusts automatically to support and reflect broader economic movements.

\paragraph{Conditional / smoothing; quality passes; $p=5.71e-02$.} Tying interest rates to the prime rate grants lenders flexibility while offering borrowers a competitive edge, as these rates move in real time with the federal funds rate and are set transparently without the delays often associated with fixed-term contractual adjustments. This method streamlines customer communication and cuts down on administrative burdens, making it easier for financial institutions to reach out to price-conscious clients efficiently. Consequently, the loan structure remains highly adaptable to changing monetary policies, ensuring that the cost of borrowing stays in sync with broader economic conditions throughout the loan's duration.
